% ============================================================
% OP_GUT2_D1_blocks_v1.tex  (EXACT INSERTION BLOCKS for review;
% NOT yet implemented). Scope per GPT verdict on recon v3:
% items 1-6 of the allowed list; short forms only; no long proof
% sketches, no D0 proof, no "all classified" language, no
% "ECT derives fermion representations", no "SM generation
% unique", no external citations.
% ============================================================

% ---------- ANCHOR 1 ----------
% Insert AFTER:
%   "Its whole-framework consistency status is recorded in the global
%    Level-4 self-consistency checklist~(\S\ref{sec:global_level4_checklist})."
% and BEFORE:
%   "\subsection{Yukawa couplings and fermion masses}"
% (end of \subsubsection*{Solution structure of the one-generation
% anomaly equations}, sec:anomaly_architecture):

\subsubsection*{P7-restricted matter audit (OP-GUT2): bounded
first pass}

\textit{Status: candidate Level~A algebra inside the restricted
class for the computed sectors; Level~B conditional as ECT
application, since the class itself is structurally imposed;
OP-GUT2 remains Open.}

\paragraph{Definition (P7-restricted admissible matter class).}
Left-handed Weyl fields valued only in tensor products of the
structures already assumed in this framework: the colour module
$E$ or $E^*$ ($\mathbf 3$ or $\bar{\mathbf 3}$), the weak
doublet/singlet structure, and a compact-$U(1)$ charge-lattice
label.
This is a structural restriction defining a test class, not a
derivation of matter representations; the exclusion of higher
colour tensors ($\mathbf 6,\mathbf 8,\dots$) and higher weak
isospin is a minimality/available-module assumption, not a
consequence of anomaly cancellation or of condensate dynamics.

\paragraph{Minimality convention.}
At most one copy of each chiral type $(R_c,R_w,Y)$; vector-like
pairs are not counted as part of the minimal chiral core; no
sterile singlets unless an explicit $\nu^c$ branch is analysed.
This is a classification convention, not a dynamical selection
rule.

\paragraph{Baseline SM-pattern sector.}
For
$\mathcal P_{\rm SM}=
\{(\mathbf 3,\mathbf 2)_{Y_Q},(\bar{\mathbf 3},\mathbf 1)_{Y_u},
(\bar{\mathbf 3},\mathbf 1)_{Y_d},(\mathbf 1,\mathbf 2)_{Y_L},
(\mathbf 1,\mathbf 1)_{Y_e}\}$
the cubic colour condition closes as
$2A(\mathbf 3)+A(\bar{\mathbf 3})+A(\bar{\mathbf 3})=2-1-1=0$,
the colour-weighted Witten count is $3+1=4$ (even), and the
remaining hypercharge conditions are exactly the one-generation
system whose solution structure is recorded
above~\eqref{eq:anom_solution_structure}.
The $\nu^c$ branches are: (a)~no $\nu^c$ (baseline);
(b)~hypercharge-neutral $\nu^c$, $Y_{\nu^c}=0$, leaving all
conditions unchanged; (c)~free-hypercharge $\nu^c$, opening the
$B-L$ flat direction recorded there.

\paragraph{Small-sector audit.}
Vector-like pairs contribute zero to all five local anomaly
conditions and an even number to the Witten count, so anomaly
conditions do not constrain vector-like additions.
Pure-singlet sets with up to four nonzero charges admit no chiral
solution: after removing vector-like $(y,-y)$ pairs and neutral
singlets, $\sum_i y_i=0$ and $\sum_i y_i^3=0$ force zeros or
$\pm$ pairs.
With no weak doublets and at most two nonzero coloured
weak-singlet multiplets, the conditions force a single vector-like
pair (and $(\mathbf 3,\mathbf 1)_0\oplus(\bar{\mathbf 3},\mathbf 1)_0$
is vector-like under colour already); the colourless doublet
sector likewise closes at minimal type-count.
With no weak doublets, $[SU(3)]^3$ requires equal numbers of
$\mathbf 3$ and $\bar{\mathbf 3}$ multiplets, excluding all odd
coloured multiplet counts.
A four-multiplet coloured weak-singlet sector, however, already
gives an explicit non-SM chiral anomaly-free family,
\[
\{(\mathbf 3,\mathbf 1)_{a},(\mathbf 3,\mathbf 1)_{-a},
(\bar{\mathbf 3},\mathbf 1)_{b},(\bar{\mathbf 3},\mathbf 1)_{-b}\},
\qquad a,b\neq0,\quad b\neq\pm a ,
\]
anomaly-free by direct substitution and containing no
gauge-invariant mass bilinear (invariance of
$\mathbf 3_{y}\otimes\bar{\mathbf 3}_{y'}$ requires $y+y'=0$,
which is excluded): $U(1)$-protected, $SU(2)$-singlet coloured
chiral quartets.
For example, $a=1$, $b=2$ is an explicit integer-charge witness.
Larger chiral patterns remain open; this is a bounded first-pass
audit, not a classification of all P7-restricted matter contents.

\paragraph{No-overclaim lemma (with witness).}
Anomaly cancellation, Witten parity, and lattice quantisation do
not by themselves select the Standard-Model matter content: the
restricted class demonstrably admits non-SM chiral anomaly-free
content (the quartet family above) alongside the SM-pattern
solution branches.
Within the class, the SM pattern is distinguished by the
minimality convention and by carrying the full
colour{+}weak{+}hypercharge chiral structure, not by anomaly
consistency alone; its dynamical derivation remains the open
problem OP-GUT2.

% ---------- ANCHOR 2 ----------
% In the closure table, REPLACE:
%   "Matter representations (OP-GUT2)
%      & Open
%      & Admissible matter class narrowed by $SU(3)^3$ and mixed-anomaly
%        cancellation requirements
%      & Sharpened \\"
% WITH:
%   "Matter representations (OP-GUT2)
%      & Open
%      & Bounded first-pass audit recorded (P7-restricted test class):
%        small sectors closed, explicit non-SM chiral quartet family
%        exhibited; anomalies alone do not select the matter pattern
%      & Sharpened \\"

% ---------- ANCHOR 3 ----------
% In the open-problems registry, REPLACE:
%   "OP-GUT2
%      & Derive fermion representations from condensate
%      & Open \\"
% WITH:
%   "OP-GUT2
%      & Derive fermion representations from condensate.
%        Bounded first-pass audit recorded in the P7-restricted test
%        class (\S\ref{sec:anomaly_architecture}); anomalies alone
%        shown not to select the matter pattern (explicit chiral
%        quartet family); minimal-core selection remains
%        convention-level.
%      & Open \\"
